\section{Introduzione}
Ogni anno, migliaia di italiani lasciano il Paese in cerca di migliori condizioni lavorative, salari più alti o più ampie opportunità professionali.
Tuttavia, i dati su cui si basano le decisioni politiche rimangono sorprendentemente inaffidabili.
La radice del problema è strutturale: l'Italia conta i propri emigranti attraverso l'ISTAT, tramite l'Anagrafe degli Italiani Residenti all'Estero (AIRE), un registro nazionale a cui i cittadini residenti all'estero sono legalmente tenuti a iscriversi~\cite{legge470_1988}.
Tuttavia, l'adesione è bassa, poiché l'iscrizione non comporta benefici immediati e la sua omissione non prevede sanzioni effettive.
Il risultato è un dataset con un bias sistematico di sottostima.
Questo studio indaga l'entità di tale bias.
Piuttosto che stimare la cifra reale dell'emigrazione da zero, adottiamo un approccio di \emph{statistiche speculari} (mirror statistics): confrontiamo i flussi di emigrazione pubblicati dall'ISTAT, derivati dalle iscrizioni AIRE, con i dati sull'immigrazione pubblicati indipendentemente dai paesi di destinazione.
Se l'Italia riporta $x$ emigranti verso la Germania in un dato anno, la Germania dovrebbe riportare circa $x$ nuovi residenti italiani.
Le divergenze sistematiche tra i due lati dello stesso flusso rivelano l'entità della sottostima~\cite{abel_cohen2019}.
Le implicazioni non sono solo accademiche.
Le statistiche sulle migrazioni influenzano le politiche del mercato del lavoro, gli investimenti e le decisioni strategiche.
Dietro ogni emigrante non conteggiato c'è un individuo che ha scelto, consapevolmente o meno, di rimanere invisibile allo Stato italiano: qualcuno che non si è registrato all'AIRE perché la procedura sembrava superflua, gravosa o semplicemente sconosciuta.
Questo comportamento collettivo distorce i dati su cui fanno affidamento i decisori politici.
Se la popolazione emigrante è sistematicamente sottostimata, le decisioni vengono prese su premesse errate.
Comprendere dove i dati falliscono è un prerequisito per correggerli.
Il resto del rapporto è organizzato come segue.
La Sezione~\ref{sec:background} fornisce il contesto sull'AIRE e i suoi limiti noti.
La Sezione~\ref{sec:data} descrive le fonti dei dati e le fasi di pre-elaborazione.
La Sezione~\ref{sec:methodology} formalizza il quadro statistico e i test d'ipotesi.
La Sezione~\ref{sec:results} presenta i risultati.
La Sezione~\ref{sec:discussion} discute le loro implicazioni con attenzione ai fattori umani che guidano la sottostima.
La Sezione~\ref{sec:conclusion} conclude il lavoro.
